\section{What the Design Bought and What It Cost}
\label{sec:discussion-synthesis}

A system paper owes the reader two things the evaluation alone does not provide:
which design choices caused each failure class, and whether a different choice
would have avoided the failure or merely relocated it. This section traces each
of the four commitments of Section~\ref{sec:design} back to its evaluation
outcome, names the alternative the designer faced at each decision point, and
states what the alternative would have cost instead. The goal is a map a tool
builder in the same design space can navigate without rerunning the experiment:
here is where each path leads, here is what we found at the end of ours, and here
is the evidence for turning earlier.

\subsection{Three distinct safety properties}

The evaluation surfaced three properties a developer needs from a tool that
rewrites files on their behalf, and they require different kinds of guarantee.

\begin{enumerate}
\item \textbf{Nothing is removed.} The store is append-only; every prior state
  is reachable by identifier. The append-only store guarantees this
  structurally, and all 10{,}237 operations in the harness preserved it.
\item \textbf{The tool operates on what you named.} The revert removes the
  feature the developer pointed at, the restore brings back the thing they
  asked for, and neither touches anything else. Three findings in the
  robustness harness broke this property: an undo that silently dropped a
  \emph{different} edit, a revert that printed success about an operation it
  had not performed, and a restore that reported ``changed nothing'' while
  emptying a file. All were recoverable---property~1 held throughout---and
  all were wrong in the direction hardest to detect, because the command's own
  report confirmed what the developer expected.
\item \textbf{New work can still be recorded.} A revert that succeeds must
  leave the repository in a state where the next save is accepted. One
  sequence in the robustness harness reached a state where every write verb
  was refused, including \cmd{undo}, because a layout record outlived the
  entity it described and the rebuild's composed image disagreed with the
  file by one byte. The developer's work was intact and unreachable.
\end{enumerate}

Property~1 is architectural and held universally. Properties~2 and~3 are
behavioural and failed on sequences involving layout records (the gap text
and ordering facts that sit between functions in a file). The append-only store
makes durability cheap, so legibility becomes the remaining problem.

The distinction matters because the three properties fail with different
visibility and different recovery costs. Property~1 failures are catastrophic
and detectable: the developer's code is gone, and they know it. Property~3
failures are non-catastrophic and detectable: the developer's code is present
but the tool refuses to record new work, and the refusal is loud. Property~2
failures are non-catastrophic and \emph{undetectable without an external oracle}:
the developer's code is present, the tool claims it did the right thing, and the
claim is wrong. In every one of eight self-reporting failures we found across the
evaluation, the data was present and the tool's account of how to reach it was
wrong. The tool exhibited what Bainbridge called the fundamental irony: the
humans most in need of monitoring are the ones least able to provide it, and here
the tool most in need of checking is the one whose reports are the only thing the
developer has to check against~\cite{bainbridge1983}. Sarter and Woods identified
the mechanism by which this gap widens: an operator whose feedback confirms a
state the system is not in stops monitoring, and every subsequent action builds on
the wrong model until a later failure exposes the divergence at a point far removed
from its cause~\cite{sarter1995}. The green checkmark over corrupted output is
exactly this structure---the developer's model (``the revert worked'') diverges from
reality (``three function headers are glued onto one line''), the divergence is
invisible because the feedback confirmed the false state, and it surfaces only when
a test runner or a colleague finds the damage.

The pattern is not unique to \sgt{}. During this evaluation, a git merge
resolution silently reverted six previously-landed fixes. Git reported the merge
as successful, the working tree was clean, and no test failed---because the same
resolution removed the tests that would have caught the reversion. The code was
present in history and absent from the working tree, with no indication that
anything had been lost. Two of the six were found only by comparing the names
in the test roster against the names at the prior commit, a check nobody was
running. This is property~2 failure in the tool the paper's baseline condition
uses, and it manifests the same way: the system reports success while silently
doing less than the developer intended.

The self-reporting pattern has a mirror that the evaluation found independently.
Seventeen files are reported as ``unrebuildable,'' but fourteen of them are files
\sgt{} decided never to record at function level---dot-paths and generated
configuration that the tier filter excluded. A developer chasing seventeen
reported breakages finds three. Over-reporting and under-reporting are the same
defect for the same structural reason: the tool's self-report confuses what it
\emph{chose not to record} with what it \emph{failed to record}, and a developer
who trusts either direction---whether the news is good or bad---acts on a model
the tool does not support.

The design implication generalises to any system that maintains a semantic record
alongside a repository. Making the store append-only gives property~1 for free
and makes property~3 failures recoverable (the developer can always fall back to
the raw files). Property~2 requires a different kind of assurance: the tool's
self-report must agree with an independently checkable outcome, and the only
oracle that can verify agreement is a probe that recomputes the result without
consulting the tool's own state. The robustness harness is exactly this probe---it
applies the operation, then checks the outcome against a reference computed from
git---and its 50\% per-session hit rate (Section~\ref{sec:session-rate}) is the
measure of how far property~2 remains from holding. Testing durability (can the
bytes be recovered?) answers a different question from testing agreement (did the
tool do what it said?), and a test suite that checks only durability passes with
flying colours on a system whose self-reports are wrong half the time.

Casserini et al.\ named this accumulation \emph{agentic entropy}: the systemic
drift between what an agent does and what the developer intended, compounding over
time and invisible to diff-based review~\cite{casserini2026entropy}. Their
diagnosis is correct---the diff tells you what changed in one commit, not how twenty
commits diverge from what was asked for---but their proposed remedy (conformity
seeding, reasoning monitoring, a causal graph of agent decisions) observes the
agent's internal trace, which is available only while the agent is running. The
record \sgt{} maintains operates at the complementary layer: it persists after the
agent exits and expresses drift in terms the developer already holds (``this
function was produced by that request''), so the entropy is visible at any future
point and actionable through the same operations that produced the work. The
distinction matters because an agent's reasoning trace is both proprietary and
ephemeral---it disappears when the session ends---while the request-to-function map
is durable and belongs to the developer.

\emph{Recoverability} must be defined precisely, because each imprecise version
produced a false violation during the harness's construction. The definition that
survived calibration: after any sequence of verbs, the bytes of every prior state
are reachable through documented commands. Reaching this statement required six
successive revisions to the test oracle, each removing an assertion about the
tool's \emph{representation} (which operation-ids a set contains) that disagreed
with the tool's \emph{content} guarantee (which bytes a file holds). The revision
sequence is itself a finding: operation-set identity and content identity are
different equivalences, and every safety claim in this domain must name which one
it means. A system can satisfy one while violating the other, and the one a
developer cares about is bytes, not identifiers.

\subsection{What each design commitment cost}
\label{sec:design-cost}

We trace each of the four commitments of Section~\ref{sec:design} back to its
evaluation outcome to show where the design has already been tested and where
it has not.

\paragraph{The function as unit (Section~\ref{sec:grain}).}
Delivered entity coverage over 42\% of files, with a threshold that loses
function identity on roughly half of all renames that co-occur with a body
rewrite (Section~\ref{sec:identity}). On the corpus the decomposed
rebuild rate is 0.23, meaning that three functions in four are absent from
the composition at \cmd{HEAD}. Ninety percent of those absences come from the
fork rule withholding records, and the parsing itself succeeds, so most of the
cost at this level actually originates from the fork rule (the next commitment)
rather than from the function-as-unit decision itself. The implication for a
designer choosing a grain: the function works as a recording unit---parsing
succeeds, identity tracking is imperfect but bounded, and the coverage spans
every file tree-sitter can read---but the cost of acting on those records
depends on what other commitments the design makes downstream.

\paragraph{Set-based versions (Section~\ref{sec:sets}).}
Gave the round-trip reversibility the walkthrough demonstrates, which the pilot
participant called ``a fast, accurate first pass that leaves me a punch list.''
That reversibility breaks twice in 10{,}237 operations, both on layout targets
no developer would name. The robustness harness found 340 operations that left
the repository in an unexpected state, and 68\% of those targeted an entity the
developer would name---but the dominant violation in those cases (157 of 232) was
a layout record being dragged along by the entity operation, while the entity
logic itself worked correctly.

The set formalism fails in two distinct ways. First, layout records: layout
facts are individually addressable because they are symbols, and individually
revertible because they are records, but no mechanism keeps them consistent with
the entities they sit between. Nine distinct defects trace to this decision, and
every robustness-harness failure entered through it. The failure has a
counterintuitive property: a lost separator is latent until both definitions
flanking it are recovered, so the number of broken files \emph{rises} as the
record becomes more complete. A developer who assumes ``more history mined,
fewer broken files'' has it backwards, and the progress bar that shows mining
advancing is not a progress bar toward a working write layer.

Second, upward closure: reverting a mid-history edit originally orphaned every
later edit that depended on
it, because the closure rule excluded them from the rebuilt set. On the study
repository this demolished 13 files when the participant asked to remove one
feature---the worst single defect we found. The fix replaced blanket exclusion
with forward subtraction (splice the removed content out of later edits rather
than dropping them), which preserves the set semantics while changing what
``remove'' means from ``exclude this and everything above it'' to ``subtract this
contribution from the tip.'' The lesson for a designer using set-based versions:
closure is correct for the append-only store (nothing is lost) but too aggressive
for the working tree (a developer removing one feature does not intend to remove
everything built on top of it). The distinction maps onto fail-stop versus
fail-silent behaviour: forward subtraction is fail-stop because the developer
sees the splice and can reject it (the preview names every function that would
change), whereas blanket exclusion is fail-silent because the developer asked
to remove one feature and thirteen files disappeared without a question.
Bainbridge's irony applies directly: the developer who installed automation to
handle the decomposition is the developer least equipped to notice when the
decomposition was too aggressive, because they have not been practising the
skill of holding it themselves~\cite{bainbridge1983}.

\paragraph{The fork rule (Section~\ref{sec:fork}).}
Protected 162 actually-forked functions and blocked mutation on 27 of 28
outside repositories. The amplification factor is 18: each forked function
withdraws records covering 17.7 function keys on average, and the total
withdrawal of 2{,}864 keys is what blocks the write guards. On the self-hosted
repository, switching the rule off raises the rebuild rate on function-level
files from 0.23 to 0.83---a 3.7$\times$ gap that the pooled figure (which
includes whole-file records that trivially reconstruct) had compressed to
2$\times$.
On the corpus, the median rebuild of 0.33 means the guard requires 1.00 and
most repositories are at one third of that. The one corpus repository that
reaches 1.00 holds 293 records where the others hold 1{,}382 to 7{,}951; it
is also the only repository where \cmd{sgt save} succeeded and the only one
the robustness harness flagged nothing on. Three separately reported positive
results trace to the same unit---a repository too thinly recorded to trigger
any of the failure modes the denser stores exhibit. The rule's read benefit is
real (the pilot participant was stopped before two conflicting edits could be
silently merged), but the write cost is eighteen times larger than the rule's
description suggests.

The fork rule is a cognitive forcing function in
Bu\c{c}inca et al.'s sense~\cite{bucinca2021}: a deliberate friction that
interrupts the developer's acceptance flow and forces engagement with the
decision the tool would otherwise make silently. Bu\c{c}inca et al.\ found
that the most effective cognitive forcing functions---those that reduced
overreliance on AI recommendations the most---were also the ones participants
rated least favourably, because the subjective cost of being stopped is felt
immediately while the objective benefit (avoiding an error) is counterfactual
and invisible. The pilot reproduced this finding exactly: the fork rule
stopped the participant before two conflicting edits could be silently merged
(the safety benefit held), and the participant declined to use the session
feature afterwards (the goodwill cost was immediate). The mechanism is that a
forcing function pays for itself only when it fires on an actual conflict, and
the developer cannot distinguish a true positive (the function genuinely
diverged) from a false positive (two edits touched opposite ends of a long
function) until they have done the inspection the rule demands. In
Bu\c{c}inca et al.'s framework, the forcing function's value depends on the
base rate of errors it prevents: if most firings are false positives, the
developer learns to treat the interruption itself as noise, and the rare true
positive arrives in a context where the developer has already learned to
dismiss.

Nelson et al.\ found that the expensive part of a merge conflict is the
semantic investigation required to understand whether both sides can
coexist, not the textual resolution~\cite{nelson2019}. The fork rule moves
that investigation from the moment the merge fails to the moment the second
edit arrives, when the developer still holds both intents. This temporal shift
is the rule's strongest claim: at the moment of the second edit, the developer
still knows what both versions intended, so the investigation costs seconds
rather than the minutes Nelson et al.\ measured at merge time.
Kazemitabaar et al.\ showed the same structure in a pedagogical setting:
cognitive engagement techniques that force students to think before accepting
AI-generated code improve comprehension at the cost of completion
speed~\cite{kazemitabaar2025}. The fork rule imposes the same trade: it stops
passive acceptance of a merge and forces the developer to state which version
they want, and the pilot confirms both sides of that exchange.

Ferino et al.\ proposed a reliance-control framework for AI in software
engineering that uses the developer's \emph{level of control} as the
mechanism for identifying both over- and under-reliance~\cite{ferino2026}.
The fork rule instantiates their framework concretely: it increases the
developer's control at exactly the point where silent merging would decrease
it (two edits to one function from different intents), and does so by making
the developer state which version they want rather than accepting whatever
the system computes. The cost is what Ferino et al.\ call unnecessary control
retention---being forced to decide when both versions are clearly compatible
---and the benefit is what they call appropriate reliance calibration: the
developer knows which merges required their judgment, and which did not.

Duan et al.\ found that distrust propagates between teammates more readily
than trust does: one person's negative experience, relayed by word of mouth,
reduces a colleague's reliance even when the colleague's own experience has
been positive~\cite{duan2026}. The fork rule has been tested only with one
person, and this contagion effect suggests that a single developer's
frustration with it could suppress adoption across a team even if the rule's
safety guarantee would have helped the others. The implication for a designer
building forcing functions into developer tools: the function must fire rarely
enough that most firings are true positives, because each false positive erodes
the willingness to engage with the next interruption, and on a team that
erosion propagates socially. On the study repository (162 forked functions in
346 commits), the fork rule fires on roughly one function per two commits; we
do not know whether that rate feels like signal or like noise to a developer
who is not the system's author. Figure~\ref{fig:fork-cost} shows the
amplification visually: the 153 files the rule blocks come from just 162
forked functions, and switching the rule off recovers them but removes the
safety guarantee.

\begin{figure}[t]
  \centering
  % The fork rule's cost: how 162 forked functions block 153 files.
% One-column figure for Section 6.3 (Rebuild completeness).
\newcommand{\FigSixCaption}{The fork rule's cost on \sgt{}'s own repository,
measured over the 252 files \sgt{} decomposes into functions. With the rule,
57 files rebuild; without it, 210 do. The 153 files the rule blocks are the
transitive cost of 162 forked functions withdrawing 2{,}864 record keys through
their dependency chains. The 42 files that fail in both cases are the residual
losses of Section~\ref{sec:unreproducible}: missing separators, incomplete chains,
and the one file whose pending record our own measurement apparatus appended.
Whole-file records (104 files, not shown) rebuild at 0.97 in both conditions
and move by exactly one file.}

\newcommand{\FigSixDescription}{Two horizontal stacked bars, each representing
252 decomposed Python files. The top bar shows 57 files rebuilding and 195 not
rebuilding with the fork rule active. The bottom bar shows 210 files rebuilding
and 42 not rebuilding with the fork rule disabled. An annotation connects the
gap between the two, labelling the 153 files as the fork rule's cost and
attributing it to 162 forked functions withdrawing 2864 record keys.}

\begin{tikzpicture}[x=1mm, y=1mm, line cap=round, font=\sffamily]

% dimensions
\def\barw{70}    % bar width in mm (represents 252 files)
\def\barh{6.5}   % bar height
\def\gapy{15}    % vertical gap between bars

% scale: 1mm = 252/70 ≈ 3.6 files; or 70mm = 252 files
\pgfmathsetmacro{\scl}{70/252}  % mm per file

% --- bar 1: fork rule ON ---
\def\yOne{20}
\pgfmathsetmacro{\rebuildOn}{57*\scl}
\pgfmathsetmacro{\blockedOn}{153*\scl}
\pgfmathsetmacro{\otherOn}{42*\scl}

\node[anchor=east, inner sep=0, text=sgtInk,
      font=\sffamily\fontsize{6.4}{7.4}\selectfont] at (-2,\yOne+\barh/2)
  {fork rule \textbf{on}};

% rebuild portion
\fill[fIdx!50, rounded corners=0.5mm]
  (0,\yOne) rectangle (\rebuildOn, \yOne+\barh);
% blocked-by-fork portion
\fill[fCache!25, rounded corners=0.5mm]
  (\rebuildOn,\yOne) rectangle ({\rebuildOn+\blockedOn}, \yOne+\barh);
% other-failures portion
\fill[sgtGrey!18, rounded corners=0.5mm]
  ({\rebuildOn+\blockedOn},\yOne) rectangle (\barw, \yOne+\barh);

% labels inside bars
\node[anchor=center, inner sep=0, text=sgtPaper,
      font=\sffamily\bfseries\fontsize{6}{7}\selectfont] at (\rebuildOn/2, \yOne+\barh/2)
  {57};
\node[anchor=center, inner sep=0, text=fCache!80,
      font=\sffamily\fontsize{5.6}{6.6}\selectfont] at ({\rebuildOn+\blockedOn/2}, \yOne+\barh/2)
  {153 blocked};
\node[anchor=center, inner sep=0, text=sgtGrey,
      font=\sffamily\fontsize{5.4}{6.4}\selectfont] at ({\rebuildOn+\blockedOn+\otherOn/2}, \yOne+\barh/2)
  {42};

% rate annotation
\node[anchor=west, inner sep=0, text=sgtInk!75,
      font=\sffamily\fontsize{5.8}{6.8}\selectfont] at (\barw+1.5, \yOne+\barh/2)
  {0.23};

% --- bar 2: fork rule OFF ---
\def\yTwo{5}
\pgfmathsetmacro{\rebuildOff}{210*\scl}
\pgfmathsetmacro{\otherOff}{42*\scl}

\node[anchor=east, inner sep=0, text=sgtInk,
      font=\sffamily\fontsize{6.4}{7.4}\selectfont] at (-2,\yTwo+\barh/2)
  {fork rule \textbf{off}};

% rebuild portion
\fill[fIdx!50, rounded corners=0.5mm]
  (0,\yTwo) rectangle (\rebuildOff, \yTwo+\barh);
% other-failures portion
\fill[sgtGrey!18, rounded corners=0.5mm]
  (\rebuildOff,\yTwo) rectangle (\barw, \yTwo+\barh);

% labels inside bars
\node[anchor=center, inner sep=0, text=sgtPaper,
      font=\sffamily\bfseries\fontsize{6}{7}\selectfont] at (\rebuildOff/2, \yTwo+\barh/2)
  {210};
\node[anchor=center, inner sep=0, text=sgtGrey,
      font=\sffamily\fontsize{5.4}{6.4}\selectfont] at ({\rebuildOff+\otherOff/2}, \yTwo+\barh/2)
  {42};

% rate annotation
\node[anchor=west, inner sep=0, text=sgtInk!75,
      font=\sffamily\fontsize{5.8}{6.8}\selectfont] at (\barw+1.5, \yTwo+\barh/2)
  {0.83};

% --- connecting annotation: the cost ---
\draw[decorate, decoration={brace, amplitude=2mm, mirror},
      line width=0.4pt, color=sgtInk!60]
  (\rebuildOn, \yTwo+\barh+1.5) -- (\rebuildOff, \yTwo+\barh+1.5);
\node[anchor=south, inner sep=0, text=sgtInk!70,
      font=\sffamily\fontsize{5.4}{6.4}\selectfont]
  at ({(\rebuildOn+\rebuildOff)/2}, \yTwo+\barh+4)
  {153 files recovered $\rightarrow$};

% --- cascade: the amplification chain ---
\def\yAnn{31}
\node[anchor=west, align=left, inner sep=0, text=sgtInk!70,
      font=\sffamily\fontsize{5.8}{7.2}\selectfont] at (0, \yAnn)
  {162 forked functions\ \ $\rightarrow$\ \ 2{,}864 withdrawn record keys\ \
   $\rightarrow$\ \ 153 blocked files};

% --- legend at the bottom ---
\def\yLeg{-2}
\fill[fIdx!50, rounded corners=0.3mm] (0,\yLeg) rectangle (3,\yLeg+2.5);
\node[anchor=west, inner sep=0, text=sgtInk!80,
      font=\sffamily\fontsize{5.4}{6.4}\selectfont] at (4, \yLeg+1.25)
  {rebuilds byte-for-byte};
\fill[fCache!25, rounded corners=0.3mm] (28,\yLeg) rectangle (31,\yLeg+2.5);
\node[anchor=west, inner sep=0, text=sgtInk!80,
      font=\sffamily\fontsize{5.4}{6.4}\selectfont] at (32, \yLeg+1.25)
  {blocked by fork rule};
\fill[sgtGrey!18, rounded corners=0.3mm] (54,\yLeg) rectangle (57,\yLeg+2.5);
\node[anchor=west, inner sep=0, text=sgtInk!80,
      font=\sffamily\fontsize{5.4}{6.4}\selectfont] at (58, \yLeg+1.25)
  {other};

\end{tikzpicture}

  \caption{\FigSixCaption}
  \Description{\FigSixDescription}
  \label{fig:fork-cost}
\end{figure}

\paragraph{Inferred names (Section~\ref{sec:names}).}
Recovered cross-file groupings that git cannot express, and systematically
over-decomposed what a developer asked for by a factor of 2 to 6 at the median
(up to 13 on the longest-session repository). On
CodeNav (the one repository where the metric discriminates), F1 is 0.66 with
precision 0.68 and recall 0.64. The margin over a file-only baseline comes
entirely from cross-file grouping. The tool's vocabulary is finer than the
developer's vocabulary, which means the operations of
Section~\ref{sec:design} require naming something smaller than what was
originally asked for. The pilot participant saw five unnamed features for one
coherent piece of work and asked why the tool was decomposing what they had done
in one sitting. On the study repositories, the feature surface is a reliable
index into history (every request's answer is reachable in at most three
commands) but an unreliable decomposition of it. Fifteen percent of labels
describe a minority of their members' actual content, and reverting a 5-symbol
feature removes 13 edits, because the structural cut crosses the episode
boundary whenever a hub function participates.

Study preparation found that 9 of 24 features in each repository owned no
behavioural code at all---communities formed entirely from positional gap-bytes
that the clustering includes for structural cohesion. Three of those husks
carried the names the study's own removal requests use, so a participant following
the map would select a feature and find nothing inside it. The fix (absorb husks
into the feature that owns their parent entities) reduced the tree from 24 leaves
to 14, but we found it only because a study task routed through it. A purity
gate that rejects a domain label when the majority of the group's symbols are
scaffold or test code would fix the 15\% label inaccuracy, but we did not build
it, because tuning on evaluation data produces a number, not a finding.
Figure~\ref{fig:decomposition} shows the ratio across all four transcript
repositories. No repository reaches 1:1, and the whiskers show that the worst
case (13 features for one request) sits on the longest-session history. The
variation between repositories is itself informative. The semi-git repository
(median 2, max 13) has the lowest median because its developer explicitly saves
after each coherent piece of work; the semipy-package repository (median 6, max
10) has the highest because its sessions interleave multiple long-running
concerns with shared infrastructure. The ratio measures how far the system's
structural vocabulary diverges from the developer's intent vocabulary, and that
divergence grows with the number of shared functions a session touches---because
each shared function creates a coupling edge that the clustering algorithm
treats as evidence of membership, pulling otherwise separate requests into a
single community. The practical reading for a developer: type \cmd{sgt log}
expecting to see your requests as you typed them, and instead find each one
decomposed into two to six pieces named after implementation details rather than
purposes. The tool's vocabulary is the code's vocabulary, not the developer's,
and the gap between the two is the gap Furnas et al.\ measured in information
retrieval~\cite{furnas1987}: the words a person uses to describe what they want
and the words a system uses to index what it has diverge by a factor that
averages 4--5 even in constrained domains.

The mismatch goes deeper than naming---it is a difference in categorical
level. A developer who typed ``add price caching'' thinks at what Rosch called
the \emph{basic level}: the category that maximises within-group similarity
while minimising between-group similarity~\cite{rosch1978}. ``The price cache'' is basic; ``the caching
subsystem'' is too abstract to act on; ``the TTL helper'' is too specific to
name a purpose. The tool presents subordinate categories---features defined by
internal structural distinctions the developer never made---and asks the
developer to act at that level (select one for revert, rename one, merge two).
The practical cost is that a developer cannot name what they want in a single
selection, because no single feature in the tool's vocabulary corresponds to
the basic-level unit they decided in. The merge and regroup commands exist to
lift subordinate features back to basic level, but that repair presupposes
that the developer knows which pieces belong together---the knowledge the tool
was supposed to supply.

The cost has a corresponding benefit that the same category theory predicts.
Rosch showed that experts operate fluently at the subordinate level where novices
cannot: a bird-watcher distinguishes warblers that a non-expert calls
``birds''~\cite{rosch1978}. A developer debugging a latency regression does not
want to revert ``the price cache''---they want to revert the TTL helper
specifically, because they know the TTL is the piece that aged the data. At that
moment the subordinate vocabulary is exactly the level they need, and the
basic-level merge would force them to take out more than they intend. The
over-decomposition is simultaneously too fine for the request-level operations the
scenarios describe and exactly right for the implementation-level operations an
expert performs daily. The design tension is real: no single granularity serves
both, and the tool chose the finer one on the ground that merging two features is
a rename (one command, no data moved) while splitting one feature after the
fact is not currently possible. The asymmetry in repair cost is what decides
the default.

\begin{figure}[t]
  \centering
  % Over-decomposition: how many sgt features one human request maps to.
% One-column figure for Section 6.7 (Grouping alignment).
\newcommand{\FigSevenCaption}{One human request maps to a median of 2--6
\sgt{} features across four repositories with full transcripts. The bar shows
the median; the whisker extends to the maximum observed. A ratio of~1 would
mean each request lands in exactly one feature. The consistently above-1 ratio
means the operations of Section~\ref{sec:design} require the developer to name
a unit smaller than what they asked for.}

\newcommand{\FigSevenDescription}{A horizontal bar chart with four rows, one
per repository: CodeNav (median 4, max 12), eico (median 3, max 8),
semipy-package (median 6, max 10), and semi-git (median 2, max 13).
A vertical dashed reference line at 1 marks the ideal ratio. All four
repositories sit above the reference, with medians between 2 and 6.}

\begin{tikzpicture}[x=1mm, y=1mm, line cap=round, font=\sffamily]

% dimensions
\def\barh{5}      % bar height
\def\gapy{10}     % vertical pitch between bars
\def\scl{5.5}     % mm per unit (1 feature = 5.5mm)
\def\maxval{14}   % axis max

% --- reference line at 1 ---
\draw[dashed, sgtGrey!60, line width=0.3pt]
  ({1*\scl}, -3) -- ({1*\scl}, 4*\gapy+\barh+2);
\node[anchor=south, text=sgtGrey, font=\sffamily\fontsize{5.2}{6.2}\selectfont]
  at ({1*\scl}, 4*\gapy+\barh+2.5) {1:1};

% --- axis ---
\draw[sgtGrey!50, line width=0.3pt] (0, -3) -- ({\maxval*\scl}, -3);
\foreach \v in {0, 2, 4, 6, 8, 10, 12, 14} {
  \draw[sgtGrey!50, line width=0.3pt] ({\v*\scl}, -3) -- ({\v*\scl}, -4.5);
  \node[anchor=north, text=sgtGrey, font=\sffamily\fontsize{5}{6}\selectfont]
    at ({\v*\scl}, -5) {\v};
}
\node[anchor=north, text=sgtInk!70, font=\sffamily\fontsize{5.6}{6.6}\selectfont]
  at ({7*\scl}, -9) {sgt features per human request};

% --- data ---
% Row 4 (top): semipy-package — median 6, max 10
\def\yD{30}
\node[anchor=east, inner sep=0, text=sgtInk,
      font=\sffamily\fontsize{6}{7}\selectfont] at (-2, \yD+\barh/2)
  {\textsf{semipy-pkg}};
\fill[fRetry!45, rounded corners=0.4mm]
  (0,\yD) rectangle ({6*\scl}, \yD+\barh);
\draw[fRetry!70, line width=0.5pt]
  ({6*\scl}, \yD+\barh/2) -- ({10*\scl}, \yD+\barh/2);
\fill[fRetry!70] ({10*\scl}, \yD+\barh/2) circle (0.7mm);
\node[anchor=west, text=sgtInk!70, font=\sffamily\fontsize{5.2}{6.2}\selectfont]
  at ({10*\scl+1.5}, \yD+\barh/2) {max 10};

% Row 3: CodeNav — median 4, max 12
\def\yC{20}
\node[anchor=east, inner sep=0, text=sgtInk,
      font=\sffamily\fontsize{6}{7}\selectfont] at (-2, \yC+\barh/2)
  {\textsf{CodeNav}};
\fill[fIdx!45, rounded corners=0.4mm]
  (0,\yC) rectangle ({4*\scl}, \yC+\barh);
\draw[fIdx!70, line width=0.5pt]
  ({4*\scl}, \yC+\barh/2) -- ({12*\scl}, \yC+\barh/2);
\fill[fIdx!70] ({12*\scl}, \yC+\barh/2) circle (0.7mm);
\node[anchor=west, text=sgtInk!70, font=\sffamily\fontsize{5.2}{6.2}\selectfont]
  at ({12*\scl+1.5}, \yC+\barh/2) {max 12};

% Row 2: eico — median 3, max 8
\def\yB{10}
\node[anchor=east, inner sep=0, text=sgtInk,
      font=\sffamily\fontsize{6}{7}\selectfont] at (-2, \yB+\barh/2)
  {\textsf{eico}};
\fill[fCache!35, rounded corners=0.4mm]
  (0,\yB) rectangle ({3*\scl}, \yB+\barh);
\draw[fCache!60, line width=0.5pt]
  ({3*\scl}, \yB+\barh/2) -- ({8*\scl}, \yB+\barh/2);
\fill[fCache!60] ({8*\scl}, \yB+\barh/2) circle (0.7mm);
\node[anchor=west, text=sgtInk!70, font=\sffamily\fontsize{5.2}{6.2}\selectfont]
  at ({8*\scl+1.5}, \yB+\barh/2) {max 8};

% Row 1 (bottom): semi-git — median 2, max 13
\def\yA{0}
\node[anchor=east, inner sep=0, text=sgtInk,
      font=\sffamily\fontsize{6}{7}\selectfont] at (-2, \yA+\barh/2)
  {\textsf{semi-git}};
\fill[fRate!35, rounded corners=0.4mm]
  (0,\yA) rectangle ({2*\scl}, \yA+\barh);
\draw[fRate!60, line width=0.5pt]
  ({2*\scl}, \yA+\barh/2) -- ({13*\scl}, \yA+\barh/2);
\fill[fRate!60] ({13*\scl}, \yA+\barh/2) circle (0.7mm);
\node[anchor=west, text=sgtInk!70, font=\sffamily\fontsize{5.2}{6.2}\selectfont]
  at ({13*\scl+1.5}, \yA+\barh/2) {max 13};

% --- legend at top ---
\def\yLeg{38}
\fill[sgtGrey!30, rounded corners=0.3mm] (0,\yLeg) rectangle (3,\yLeg+2.2);
\node[anchor=west, inner sep=0, text=sgtInk!80,
      font=\sffamily\fontsize{5.2}{6.2}\selectfont] at (4, \yLeg+1.1)
  {median};
\draw[sgtGrey!60, line width=0.5pt] (18,\yLeg+1.1) -- (22,\yLeg+1.1);
\fill[sgtGrey!60] (22,\yLeg+1.1) circle (0.6mm);
\node[anchor=west, inner sep=0, text=sgtInk!80,
      font=\sffamily\fontsize{5.2}{6.2}\selectfont] at (23.5, \yLeg+1.1)
  {maximum};

\end{tikzpicture}

  \caption{\FigSevenCaption}
  \Description{\FigSevenDescription}
  \label{fig:decomposition}
\end{figure}

\paragraph{The commitments as a system.}
The four commitments do not fail independently. The function-as-unit decision
creates the records the fork rule must protect, the fork rule's amplification
blocks the write guards the set formalism requires, and the set formalism's
layout seam produces the violations the inferred names must absorb into a
coherent display. This means that fixing any one commitment's failure often
exposes or worsens another's. Loosening the fork rule (to stop blocking 27 of 28
outside repositories) gives up the consistency guarantee the set formalism
depends on. Eliminating layout records (to remove 68\% of violations) degrades
the rebuild, because the whitespace between functions is what makes the
reconstructed file compilable. Sharpening the clustering (to reduce 2--6$\times$
over-decomposition) requires a signal the fork rule suppresses, because a forked
function contributes no included version and therefore no coupling edge.

On the self-hosted repository, the full causal accounting closes to zero.
Dependency grounding excludes 7 operations. The fork rule excludes 1{,}298, and
because it drops both tips and their upward closures, those 162 forked functions
withdraw 2{,}864 function records. The pending working tree accounts for 292
more. Nothing is left over: the recorded ideal rebuilds exactly as many files as
a walk through grounded, fork-free records produces. The mechanism is not a bug;
it is three separately reasonable decisions---never silently pick a side when two
records claim the same function, require 1.0 reconstruction before writing, and
leave the test oracle unconfigured---that compose into a tool that cannot record
on 27 of 28 repositories including its own. Every fork is resolvable, and the
thing that would resolve them automatically is a test runner nobody configured.

The interaction teaches a design lesson that the per-commitment analysis
obscures. A designer in this space faces a choice not between four independent
bets but between packages of mutually-constraining decisions. The package we
chose optimises for correctness of a single write operation (a revert removes
exactly the named set and nothing else) at the cost of availability (the write
guards fire often) and learnability (the vocabulary is too fine). A different
package---coarser grain, weaker safety, simpler surface---would trade those costs
for a higher rate of silent errors at revert time, which is the cost
Section~\ref{sec:scenarios} argues developers already pay and do not notice. The
study is designed to measure whether they notice, and the answer determines which
package is the better bargain.

\subsection{The codebook that was wrong}
\label{sec:codebook-wrong}

We coded 60 sampled errors from two repositories (CodeNav and eico) against
a pre-registered four-code taxonomy: split (one intent across two features),
merge (two intents in one feature), mis-assign (correct groups, wrong
membership), and rename (correct membership, wrong label). The pre-registered
codes explained 6 of 30 errors in CodeNav and 15 of 30 in eico. On both
repositories, the dominant mechanism was a disagreement about what counts
as one request, which the clustering handled correctly.

Two mechanisms accounted for the residual. First, \textbf{consecutive turns of
one feature}: a developer who types ``resume'' and then ``ok sure'' between
substantive requests creates turn boundaries that the ground truth treats as
request boundaries, but \sgt{} groups the symbols from both turns into one
feature, which is the semantically correct answer. On CodeNav this mechanism
produced 15 of 30 coded items. Second, \textbf{continuation tokens as intent}:
on eico, 88\% of turns are contentless directives (``resume'', ``move on'',
``(a)''), and the ground truth attributes each turn's output to that turn's
text, producing a positive base rate of 47\% against a signal that carries no
semantic content. \sgt{} correctly refuses to separate work that differs only in
which continuation token preceded it.

This finding generalises. Anyone measuring intent recovery against real agentic
transcripts will face the same problem, because a per-turn ground truth does not
yield a per-turn metric on histories where most turns carry no semantic content. The fraction of contentless turns varies from 0\% to 88\%
across the four repositories we hold, and it is not predictable from session
length or repository size. A metric computed on a history with 88\%
continuation tokens measures the ground truth's granularity rather than the
system's accuracy.

The implication extends beyond our system. The standard evaluation protocol in
change untangling research constructs ground truth from developer-authored commit
messages, treating each commit as one intent~\cite{herzig2013,dias2015}. That
protocol assumes the commit boundary aligns with the intent boundary, which
Herzig and Zeller showed fails on 15\% of commits in well-maintained
projects~\cite{herzig2013}. Agent transcripts make the problem worse, not better:
a commit boundary is at least a deliberate act (the developer typed \cmd{git
commit}), whereas a turn boundary is an artefact of the chat interface (the
developer pressed enter). The pre-registered codebook assumed that the
commit/turn boundary would be wrong in predictable ways---split, merge,
mis-assign, rename---and found instead that the dominant residual was a
disagreement about what counts as a boundary at all. The ``continuation token''
mechanism is specific to agentic transcripts, but the underlying phenomenon
(the ground truth's grain is finer than the intent it claims to record) will
appear in any corpus where the recording unit (turn, commit, save) can be
triggered without semantic content. A future evaluation protocol for intent
recovery should filter contentless boundaries before computing any metric,
or else report the contentless fraction alongside the score so that a reader
can judge whether the number measures the system or the ground truth.
This result also explains why the clustering in Section~\ref{sec:names} uses
co-change within saves rather than within turns: a save is a developer-initiated
boundary that contentless turns do not trigger, so the temporal signal the
clustering reads is already filtered.

The one mechanism that \emph{is} an error in the clustering---intent sharded by
directory (F32)---appeared in both repositories and has a consistent shape: two
leaves carry the same label and the same intent, separated only because their
symbols live in different directories, with no interior node uniting them.
This is an addressable defect: a merge pass that detects identically-labelled
sibling leaves would eliminate it without changing any parameter.

\subsection{Comparison with alternatives}
\label{sec:comparison}

The evaluation data lets us compare against what existing tools offer as
well as against our own aspirations. Table~\ref{tab:comparison} summarises the
tradeoffs.

\begin{table*}[t]
\centering
\small
\begin{tabular}{p{1.8cm}p{2.1cm}p{2.1cm}p{2.1cm}p{2.1cm}p{2.1cm}}
\toprule
Property & \textbf{Darcs/Pijul} & \textbf{Jujutsu} & \textbf{GitButler} & \textbf{sem} & \textbf{\sgt{}} \\
\midrule
Unit of record & Line-level patch & Commit (change id) & File-level change & Function (derived cache) & Function (committed record) \\
Cross-file dep. & Not tracked & Not tracked & Not tracked & Reported, not enforced & Tracked and enforced \\
Selective revert & Pull/unpull a patch & Backout a change & Undo a file assignment & N/A (read-only) & Revert a named feature \\
Adoption cost & Replace git & Replace git & Add to git workflow & None (cache alongside) & Add to git workflow \\
Failure mode & Silent (line-level) & Silent (commit-level) & Silent (file-level) & N/A & Fork rule stops \\
\bottomrule
\end{tabular}
\caption{Design tradeoffs across five systems that address change grouping.
Darcs and Pijul compute dependency from lines; Jujutsu replaces git's commit
model with stable change identities; GitButler assigns files to virtual branches;
sem derives a function-level cache but takes no action on it; \sgt{} commits its
record and acts on it. Each step down in granularity enables a more precise
revert but introduces a failure mode the coarser system never encounters.}
\label{tab:comparison}
\end{table*}

Darcs computes dependency from line adjacency~\cite{roundy2005}, so pulling a
patch that edits \sym{api.py::handle} without the patch that defined
\sym{Retry} in another file produces a reference to a name that does not exist,
with no refusal and no warning. \sgt{}'s closure rule prevents this by refusing to write rather than writing a
file that does not parse, and the fork rule's amplification is the cost. The
evaluation shows that this conservatism is too expensive on repositories \sgt{}
did not record live (27 of 28 refused entirely), but the pilot shows it catching
the case it was designed for (two edits to one function from separate requests,
neither aware of the other).

sem offers the entity-level answers without the action layer~\cite{sem2026},
which makes it safe to adopt and safe to abandon. A derived cache that any
teammate's commit can silently invalidate cannot be the record an operation acts
on, so sem's placement is a deliberate choice to trade power for safety. The
evaluation suggests this may be the right stopping point for a tool that cannot
yet guarantee its write operations: a tool that reports without acting inherits
property~1 automatically (there is nothing to corrupt) and sidesteps property~3
entirely (there is no write guard to trigger). The pilot confirms the
consequence: the participant independently described sem's value proposition by
rephrasing the revert as ``a fast, accurate first pass that leaves me a punch
list.'' The apparatus pilot showed that the harder problem was reaching the
answers at all, since the primary inspect verb rejected 6 of 10 queries
(Section~\ref{sec:pilot-apparatus}). sem does not face this problem because its
commands accept whatever git already accepts (commit hashes, branch names), so a
developer's existing vocabulary works without learning a new selector namespace.
\sgt{} cannot offer the selective
revert, the restore, or the switch from sem's position, because those operations
require writing the subset back to disk---the record must be committed alongside
the code so the rebuild is deterministic. Those operations are what motivated
\sgt{}'s decision to commit its record, and the evaluation shows they are also
where every failure lives. The honest reading of this comparison is that \sgt{}
took on an obligation sem declined, and the obligation has not yet paid for
itself in the write layer.

GitButler assigns file-level changes to virtual branches~\cite{gitbutler2026},
which handles the common case where one feature per file suffices. The scenario at the centre of
this paper---three requests all editing \sym{api.py::handle}---is the case
where file-level assignment cannot express the grouping. \sgt{}'s function-level
record handles this case and pays for it with the fork rule, the layout seam,
and the identity threshold. The comparison teaches a design lesson: every step
downward in granularity (file to function, commit to edit) introduces a new
class of consistency problem (layout records, identity tracking, fork
amplification) that the coarser grain never had to solve, and the evaluation
evidence for whether that cost is justified remains thin---we measured hub
absorption (19--41 symbols per shared feature on the study repositories) but not
the frequency with which real agent sessions produce the multi-request
single-function case that motivates the finer grain.

The tradeoff is not symmetric. Moving from file to function introduces three
problems (layout consistency, identity across renames, fork amplification) but
addresses one: the tangled function that holds edits from multiple requests. If
that case is rare, the three problems dominate and the finer grain is a net loss.
If it is common---and agent sessions make it common, because a model asked to
``add caching'' and ``add rate limiting'' will change the same handler
function---the finer grain addresses the single case that the coarser grain
cannot express. The question a tool builder faces is whether to optimise for
the common case (most files serve one feature; file-level grouping suffices) or
for the hard case (one function serves multiple features; only function-level
grouping can separate them). GitButler chose the common case and pays nothing
for the hard case because it does not attempt it. \sgt{} chose the hard case and
pays for it on every file, including the majority where the hard case does not
arise. A tool builder choosing a granularity should treat our failure catalogue
as a lower bound on the cost of going finer, and GitButler's adoption as
evidence that the coarser grain satisfies most workflows today.

Jujutsu keeps git's line-level recording and adds what git lacks: a change
identity that survives rewriting, divergence detection when two worktrees evolve
the same change differently, and a working-copy model that lets a developer
edit freely with no staging ceremony~\cite{jujutsu2026}. These are exactly the
properties \sgt{} provides at the function level, delivered at the commit level
with no parsing, no clustering, and no language dependency. The cost is that
Jujutsu's change is still a commit: it cannot express ``the caching edits inside
\sym{api.py::handle}'' as a unit, because it cannot see inside a file. The
developer who wants to take the caching out and leave the retry policy intact
faces the same hunk-selection problem they face in git, in a cleaner interface
but with the same decomposition burden. A fair reading is that Jujutsu solved the
workflow problems git's model created (detached HEAD, lost rebases, opaque
staging) while \sgt{} attempted the representation problem underneath them (the
commit is too coarse for a request), and the two efforts address overlapping
symptoms from different root causes. A developer who adopts Jujutsu still needs
to supply the decomposition at revert time; a developer who adopts \sgt{} still
needs Jujutsu's workflow improvements for the commits \sgt{} cannot decompose
(configuration, prose, generated files). The two are complementary rather than
competing, and we have not tested whether \sgt{} works atop Jujutsu's storage
model.

The untangling literature (Herzig and Zeller~\cite{herzig2013}, Barnett et
al.~\cite{barnett2015}, Dias et al.~\cite{dias2015}, Shen et
al.~\cite{shen2021}) addresses the same problem after the fact: given a tangled
commit, recover the partition. Their evaluation compares a computed partition
against a human-labelled ground truth, and Section~\ref{sec:codebook-wrong}
reports what we found when we tried the same approach on agent transcripts: the
ground truth's grain disagrees with the semantic grain so thoroughly that the
metric measures the disagreement rather than the system. \sgt{} differs from this
literature in timing rather than technique---the coupling signal and community
detection we use are the same signals they use---and the timing difference
matters because the record is written while the session is still running, before
the developer has lost the words that name each request. Dias et al.'s design
comes closest to ours: they record fine-grained IDE events as the developer types
and cluster them into ``development sessions''~\cite{dias2015}. They stopped at
the boundary of the editor, as every tool in Section~\ref{sec:related} did, and
the reason none of them crossed it is the same reason we argue the crossing is
now justified: a record that sits outside the repository cannot be the record an
operation acts on, and until a developer needed operations they could not already
perform by hand, the crossing was cost without benefit.

The comparison above treats the tools as interchangeable alternatives. They are
not. The shift from inline completion to CLI agents changes the level at which
version control must operate. Murphy-Hill et al.'s field study found that
CLI-based coding agents produce a sustained 24\% lift in merged pull requests
over four months~\cite{murphyhill2026}, while cursor-style IDE completion shows
velocity gains that dissipate within two months as complexity accumulates and
developers spend their gains on rework~\cite{he2026}. The persistence difference
is structural: an inline completion changes one expression inside a function the
developer is already reading; a CLI agent session changes fourteen functions
across five files in response to a sentence the developer never read back. The
first is a line-level event that git already records at the right grain. The
second is a session-level event whose decomposition no existing tool preserves,
and it is the second that is growing. Among professionals using AI daily, the
share who primarily direct an agent (as opposed to accepting completions) rose
from 12\% to 38\% in the twelve months preceding the survey that measured the
trust decline~\cite{stackoverflow2025}. The tangled-function scenario that
justifies function-level recording---three requests editing one handler---is
produced specifically by the session mode, because a completion does not cross
file boundaries and an agent does. The design space this paper occupies is
therefore not a permanent fixture of development; it is the space that opens when
the dominant interaction mode shifts from completing a line the developer started
to executing a request the developer specified. If that shift reverses, the
decomposition problem reverses with it and the finer grain becomes pure cost.

Three recent empirical studies suggest the shift is accelerating rather than
reversing, and that its artefacts persist. Rahman and Shihab tracked 200{,}000 code
units across 201 open-source projects and found that agent-authored code survives
\emph{longer} than human code---a 15.8 percentage-point lower modification rate---while
carrying a modestly elevated corrective-change rate (26.3\% vs.\
23.0\%)~\cite{rahman2026survive}. Code that persists longer while accumulating more bug
fixes is code whose provenance question becomes more acute over time, not less: each
surviving function is a future undo that requires tracing back to the session that
introduced it. Liu et al.\ confirmed the scale of what survives: across 302{,}600
AI-authored commits in 6{,}299 repositories, 22.7\% of the issues the AI
introduced---code smells, correctness defects, security vulnerabilities---still exist at
the latest version~\cite{liu2026debt}. Nearly one in four problems becomes permanent
technical debt, persisting long after the session that produced it has ended and the
developer has forgotten what they asked for. Sawada et al.\ closed the triangle from the
maintenance side: AI-generated files receive less frequent maintenance than human-authored
code, and when maintenance does occur, human developers perform the large majority of
it~\cite{sawada2026maintenance}. The pattern across all three is that AI-generated code
enters the repository, persists, accumulates corrections performed by people who did not
write it, and the longer it persists without provenance, the harder each correction
becomes---because the developer maintaining it cannot distinguish what was asked for from
what was produced incidentally.

\subsection{When not to adopt this design}
\label{sec:non-adoption}

The conditions under which this design should not be adopted follow directly from
the comparison.  A repository whose functions serve one purpose each---the common
case file-level assignment handles---gains nothing from function-level recording
and pays its full cost (layout consistency, identity tracking, fork
amplification).  A team whose members edit the same functions daily gains the
fork rule's safety at a queue length that erodes goodwill faster than it prevents
errors (Section~\ref{sec:forkoff}).  A repository too thinly recorded for any
failure mode to trigger---the Complex-YOLOv3 case, 293 records where others hold
thousands---pays nothing but also learns nothing; git suffices because the tool's
guards never fire.  And any repository whose primary concern is workflow
(staging ceremony, detached HEAD, lost rebases) is better served by Jujutsu,
which solves those problems without parsing a single file.  The design is
justified only when function-level tangling---multiple requests building
overlapping logic inside shared functions---is common enough to outweigh the
three problems the finer grain introduces.  The criterion is the interaction
mode, not the presence of AI tooling.  A team that primarily uses inline
completion---accepting or rejecting one expression at a time inside a function
they are already reading---produces the same edit structure it produced before,
and git records it at the right grain.  A team that primarily directs CLI agents
in session mode produces the tangled-function structure this design addresses,
and the shift between these modes is ongoing: the share of developers who
primarily direct an agent rather than accept completions tripled in twelve
months~\cite{stackoverflow2025}.  Handa et al.\ confirmed this split from the
opposite direction---observed conversations rather than self-report---and found
that 43\% of AI usage across four million interactions constitutes automation
(the user specifies intent and receives completed work with minimal
involvement), while 57\% constitutes augmentation (the user iterates, learns,
and stays in the loop)~\cite{handa2025tasks}. The convergence from two methods
(38\% self-reported agent direction, 43\% observed automation) on roughly the
same figure suggests that the mode \sgt{} addresses is not a fringe behaviour
but the minority that is also the fastest-growing segment.  The evaluation
cannot yet say how common function-level tangling is on teams larger than one,
but the mode that produces it is growing. Nanda et al.\ provided a base rate for the provenance problem:
30\% of production agent trajectories exhibit specification drift---the agent
deviates from what the developer asked for---and even the 70\% that complete
without incident still interleave edits from multiple intents inside a single
session~\cite{nanda2026wink}. Their Wink system corrects drift in real time, but a
specification drift that Wink does not catch, or that the developer does not notice
until a week later, becomes a provenance error that only surfaces when someone
tries to undo one request without disturbing the others. The 30\% figure bounds
how often a session leaves the repository with code filed under the wrong request,
and at that rate a developer whose week involves ten agent sessions carries three
provenance errors into the following week by default.

\subsection{Design implication: build the read layer before the write layer}

The evaluation, the pilot, and the comparison converge on one recommendation for
a designer facing the same problem: build the read layer first, ship it alone,
and defer the write operations until the developer trusts the reports.

The urgency of this recommendation comes from the wider ecosystem. Developers are
adopting AI tools nearly universally (84\% report using them) while trusting
their output less each year---accuracy trust fell from 40\% to 29\% in a single
survey cycle~\cite{stackoverflow2025}. The gap between adoption and trust means
that a tool entering this space inherits a context in which developers already
expect automation to produce output they must verify, and already expect that
verification to be expensive. A write layer whose first encounter produces a
violation lands in an environment pre-loaded with scepticism, while a read layer
that shows correct attributions earns trust against a low baseline---precisely
because the developer expected it to be wrong and found it was right. He et
al.\ confirmed the downstream half of this dynamic: velocity gains from
AI-assisted development dissipate within two months on repositories where
complexity rose~\cite{he2026}. Speed without legibility is self-defeating, and a
read layer is legibility. Faros AI's engineering telemetry across 22{,}000
developers confirmed the pattern at fleet scale: under high AI adoption,
throughput rose (epics per developer up 66\%, PRs merged up 16\%) while quality
degraded faster---bugs per PR up 29\%, median review time up
5$\times$, production incidents per PR up 3$\times$, and code churn up
10$\times$~\cite{faros2026whiplash}. The throughput gains are real; the quality
deficits are what happens when no mechanism keeps comprehension proportional to
output, and the review-time explosion is what happens when the burden of
restoring that proportion falls on human reviewers reading code nobody
understands. A read layer that maintains the request-to-function map is
infrastructure against exactly this dynamic: it gives the reviewer a vocabulary
for what the code is doing there, before they spend five times longer working it
out from the diff alone. Lyu et al.\ gave the accumulated cost a name:
\emph{comprehension debt}---code outpacing understanding---which their
practitioner interviews surfaced as one of six central challenges in building SE
agents~\cite{lyu2026practitioners}. Storey's Triple Debt Model distinguishes three kinds of deficit that degrade
a codebase: technical debt (in the code), cognitive debt (in the team's shared
understanding), and \emph{intent debt}---the absence of externalized rationale
that developers and agents need to work safely with
code~\cite{storey2026tripledebt}. AI-generated code can be technically clean
while carrying maximum intent debt: every function works, and nobody knows
which request produced it or why. The read layer prevents intent debt from
accruing by externalizing the rationale in real time: it records which request
produced which function, so that the ``why'' survives the session boundary. It
also slows cognitive debt, because the developer who reads the attribution is
maintaining the shared understanding that would otherwise erode through disuse.
Comprehension debt, in Lyu et al.'s terms, accrues whenever a session
ends with code nobody has mapped back to the request that produced it, and the
read layer is the mechanism that prevents that accrual by maintaining the map in
real time. It is the thing that makes the speed worth keeping.

Vella and Blincoe measured the cost of speed without comprehension longitudinally:
across 95 matched developers surveyed six months apart, productivity perceptions
held stable (84\% reporting improvement at both time points) while the share
reporting worsened developer experience nearly doubled from 14\% to
27\%~\cite{vella2026longitudinal}. They named the new category of work this shift
produces \emph{supervisory engineering}---direction, evaluation, and correction of
AI output---and found that flow state and cognitive load both eroded over the six
months while feedback loops improved. The paradox is precise: developers feel
productive while losing the experience dimensions (flow, cognitive ease) that
depend on understanding what is happening in their code. A read layer that
maintains the request-to-function map is infrastructure for supervisory
engineering---it gives the evaluation step (``what did the agent produce?'')
a concrete answer---and without it, the supervisory work degenerates into the
manual recovery the five scenarios describe.

The calibration failure can be worse than Vella's paradox suggests. Becker et
al.\ ran a randomised controlled trial on sixteen experienced open-source
developers working on their own repositories---projects they had maintained for
years---and found that AI coding tools made them 19\% slower on average, while
the same developers perceived a 20--24\% speedup~\cite{becker2025metr}. The
gap is not mere optimism; it is an inversion: the developers' self-report
pointed in the opposite direction from the measurement. Rozenblit and Keil's
illusion of explanatory depth (Section~\ref{sec:related}) supplies the
mechanism: the developers understood their sessions at the purpose level (``I
asked for X and got working code'') and mistook that understanding for a
productivity gain, while the mechanism level---the one that determines actual
speed---was absent from their self-model entirely. A developer who cannot
perceive the productivity effect of a tool certainly cannot perceive the
comprehension gap that same tool opens---the accumulating distance between what
is in the repository and what the developer understands about it. External
ground truth is what corrects the calibration, and a read layer that shows which
request produced which function is a form of ground truth the developer can
check against their own understanding. If the attribution surprises them, the
gap is already there and growing.

Parasuraman's model explains the mechanism. Automating stages one and two (parsing,
clustering, dependency display) while leaving stage three to the developer
(which feature to name, which revert to confirm) is both safer and independently
useful~\cite{parasuraman2000}. The pilot confirmed this empirically: the
participant adopted the read layer on first contact and declined the write layer
after one session. The corpus confirmed it structurally: on 27 of 28 outside
repositories, the write layer had nothing to offer because its preconditions
were unmet, while the read layer (attribution, dependency display, feature
listing) ran on all 35.

The practical consequence is that a tool in this design space should ship its
read layer without waiting for write-path stability. The argument has five
parts, each grounded in a distinct mechanism.

First, \emph{commitment timing}. Thimbleby argued that good design delays
commitment until the user holds the information the commitment
requires~\cite{thimbleby1988}. A developer who has not yet seen the tool's
attributions does not hold the information needed to decide whether to trust
its reverts; the read layer is what supplies that information. Shipping the
write layer first asks the developer to commit (accept a revert, resolve a
fork) before they have any basis for judging whether the tool is right.
Sarkar and Drosos confirmed this pattern empirically in extended vibe coding
sessions: trust develops through iterative verification rather than blanket
acceptance, and expertise redistributes toward rapid code
evaluation~\cite{sarkar2025}. The read layer is the substrate on which that
iterative verification operates---a developer who has seen ten correct
attributions has built the basis for judging the eleventh, which is the
mechanism Thimbleby's principle requires.

Second, \emph{reliance decomposition}. Schemmer et al.\ decomposed reliance into
two independent measures: the rate at which a user follows the system when it is
right (RAIR), and the rate at which they override it when it is wrong
(RSR)~\cite{schemmer2022}. A read layer builds RSR directly, because a developer
who has seen the attribution can judge whether a proposed revert targets the
right set, and override it when it does not. A write layer shipped without that
foundation offers no basis for override, and a single visible failure then
collapses both metrics---the developer stops following even correct
recommendations because they have no independent basis for distinguishing them.
Section~\ref{sec:session-rate} shows the 50\% per-session violation rate; a tool
whose write layer fails in half of all sessions cannot build RAIR faster than
its own failures erode it. Dixon and Wickens showed that the damage mechanism
is specific to signal type: misses (the system fails to alert) erode reliance,
while false alarms erode compliance~\cite{dixon2007}. The write layer's
self-reporting failures are misses---the system should signal a problem and
instead confirms success---so they damage precisely the reliance the read layer
needs to have already built. A read layer that works correctly for twenty
minutes builds reliance through demonstrated accuracy; a write failure in minute
twenty-one damages that same channel, and the asymmetry (one miss undoes many
correct silences) is what makes temporal ordering matter.
Tang et al.\ confirmed the pattern at ecosystem scale: across 20{,}574
real-world coding-agent sessions, inaccurate self-reporting accounts for
22.6\% of all developer--agent misalignment episodes, and its share is
\emph{rising} even as overall misalignment rates
decline~\cite{tang2026agentfailures}. The technical failures agents learn to
avoid first---wrong implementation, wrong diagnosis---are being replaced by
interaction-level failures that look like successes, because a model that
produces correct code more often still reports completion when it has not
achieved what the developer asked. The implication for temporal ordering is
direct: the failure class the read layer must survive is not shrinking; it is
becoming the dominant residual, and a tool that ships its write layer into that
environment inherits a rising base rate of undetected misreporting against which
it must build reliance from scratch.

Third, \emph{situation awareness}. Endsley and Kiris showed that operators
relegated to supervisory monitoring under automation lose the active engagement
needed to maintain awareness of system state, and respond more slowly and less
effectively when anomalies occur~\cite{endsley1995}. A developer who delegates
code generation to an agent is in precisely this monitoring posture: they specify
intent and observe output, but do not actively engage with the line-by-line
decisions the model makes. Bainbridge's irony follows: the skill of holding the
decomposition degrades through disuse, and the degraded skill is exactly what the
developer needs when the automation fails~\cite{bainbridge1983}. Risko and
Gilbert's cognitive offloading framework describes the stronger case: people
strategically delegate cognitive work to external resources when internal
processing is costly, and the offloaded information is never encoded in memory
at all~\cite{risko2016}. A developer who offloads production to an agent has
not forgotten the decomposition; they never held it, because the offloading
\emph{succeeded}. Situation awareness support for this developer must compensate
for absent encoding, not for degraded recall. The read layer
(attribution, dependency display, consequence preview) functions as what Endsley
called a ``situation awareness support system''---it restores the developer's
model of which functions serve which purpose, \emph{without} requiring that they
build it from reading the code. Without that support, the developer who needs to
intervene on a failed revert is in the out-of-the-loop position: they have lost
both the comprehension the automation took from them and the active engagement
that would have maintained it. Cao et al.\ confirmed the dependency
experimentally: explanations improve appropriate reliance on automation only when
they align with the user's prior intuition~\cite{cao2024}---a preview that names
consequences is useful precisely because the read layer has already built the
intuition it appeals to. Without that prior comprehension, the same preview is an
assertion the developer cannot evaluate.
The read layer therefore serves a second function beyond trust building: it
counteracts the deskilling the agent itself produces. Shen and Tamkin measured
the cost directly: in a controlled experiment, developers who used AI assistance
scored 17\% lower on post-task assessments than those who coded by hand (Cohen's
$d$=0.74), with debugging---the ability to identify when code is wrong and
why---showing the largest degradation~\cite{shen2026skills}. The skill most damaged
is precisely the one a developer needs to evaluate a revert preview: judging
whether the consequence list is complete requires understanding what each function
does well enough to notice what is missing. Bastani et al.\ found the same pattern
in a learning context: developers who practiced with a generative model performed
worse on subsequent assessments without it~\cite{bastani2025}. Catalan et al.\
observed the mechanism in real time: cognitive engagement with agent output declines
as tasks progress, and current agentic assistants provide no affordance for the
reflection that would maintain it~\cite{catalan2026notreading}. The developer stops
reading before the session ends---not because the code is bad, but because the
interface offers no way to absorb provenance without reading every line.

Shen and Tamkin's archetypes make the design implication specific. Their
highest-scoring group (``conceptual inquiry'') asked only conceptual questions and
coded independently---they used the AI to understand rather than to produce. Their
lowest-scoring groups delegated production entirely or drifted from independence
into delegation as the session progressed. A read layer is infrastructure for the
conceptual-inquiry pattern: it answers ``what did the agent produce?'' and ``which
request is this function from?'' without producing any code itself, so a developer
who consults it is exercising precisely the engagement pattern that preserved skill
in the experiment. A write layer that ``just works'' enables the delegation pattern
that degraded it. The temporal strategy---read first, write later---is therefore a
skill-preservation sequence as much as a trust-building one: it keeps the
developer in the evaluation posture that Shen and Tamkin showed protects learning,
while deferring the automation that their progressive-reliance archetype showed
erodes it.

A developer who regularly reads \sgt{}'s attribution is exercising the
comprehension skill in miniature: they see the grouping, compare it against what
they believe, and correct it when it diverges. The attribution is the affordance
Catalan et al.\ found missing---a structure that sustains engagement by presenting
provenance as a scannable list rather than a wall of code. That exercise maintains
the comprehension that agent-mediated development otherwise erodes, and the
comprehension so maintained is what makes the write layer's previews evaluable
when the developer eventually encounters them.

Fourth, \emph{forcing-function calibration}. Bu\c{c}inca et al.\ found that the
cognitive forcing functions which reduced overreliance most were also the ones
participants rated least favourably~\cite{bucinca2021}. A forcing function
encountered before the user has built a model of the system's accuracy is
experienced as pure cost, because the user cannot tell whether the interruption
prevented an error or merely wasted time. The fork rule is a forcing function
(Section~\ref{sec:design-cost}), and the pilot confirms that it was experienced as
pure cost: the participant who encountered it had no basis for knowing that the
interruption prevented a real conflict, because they had not yet seen enough
correct attributions to calibrate. A read-only release whose attribution is
correct builds both the trust and the comprehension that a later write release
depends on---including the comprehension that makes a forcing function feel like
signal rather than noise.

Fifth, \emph{mode compatibility}. Barke et al.\ found that developers interact
with AI code tools in two distinct modes: acceleration (the developer knows what
they want and uses the tool to get there faster) and exploration (the developer
does not know what they want and uses the tool to discover
options)~\cite{barke2023}. The read layer supports both modes: in acceleration,
attribution confirms the developer's existing model; in exploration, attribution
reveals structure the developer has not built yet. The write layer supports only
acceleration, because a developer in exploration mode cannot verify whether a
revert removed the right set---they lack the model against which to check.
Mozannar et al.\ modelled this verification cost explicitly: how long a
developer spends evaluating a code suggestion depends on their existing
understanding of the code at the moment the suggestion
appears~\cite{mozannar2024}. A developer with no prior model of which functions
serve which purpose faces a verification cost that exceeds the cost of doing the
revert by hand, at which point the tool's contribution is negative. The read
layer is what builds that prior model; without it, the write layer's
verification cost is unbounded. Xu et al.\ reached the same conclusion from
the opposite direction: they argued that intelligent coding systems should
produce \emph{justifications} alongside code---explanations that bridge model
reasoning and user understanding---and identified cognitive alignment (matching
how the user thinks about the task) as the critical property of such
justifications~\cite{xu2025justifications}. Attribution \emph{is} a
justification in their sense: it says which request produced which code, in the
vocabulary the developer used to ask for it, so the developer can evaluate the
record against what they remember rather than against what the code says.
Velasco et al.\ posed the question at the model level---tracing generated code
back through prompt, training data, and internal
representations~\cite{velasco2026provenance}---and that explanation is orthogonal
to \sgt{}'s: one tells the developer why a line has these particular tokens
(the model's knowledge), the other tells the developer why this function is in
their project at all (their own request). The second is what makes a revert
evaluable, because ``I asked for a price cache'' is a fact the developer already
holds and can check the attribution against.

These five mechanisms converge on a prediction Rieman's field study of
exploratory learning confirms empirically~\cite{rieman1996}. Users encountering
new software try low-cost operations before high-cost ones, explore when a real
task is at hand, and adopt features whose failures they can detect and reverse
before features whose failures propagate silently. A read-only layer is exactly
what this adoption trajectory predicts: a developer on first contact inspects
the attribution, checks whether it matches what they believe, and either trusts
the label or renames it---all at zero risk to their files. Only after several
such confirmations do they attempt an operation that modifies state. Our pilot
replicated this trajectory in miniature: the participant used \cmd{sgt show}
and \cmd{sgt log} within the first three minutes, attempted \cmd{sgt revert}
only after twenty minutes of reading, and articulated the exploratory strategy
explicitly: ``I trust the reports but I don't trust the operations yet.'' A
tool that ships only its write layer forces the developer past the exploration
phase into commitment on first contact, which is the structure Rieman's subjects
consistently refused to adopt. Wang et al.\ framed the same trajectory as a
research agenda: coding agents should be evaluated on task alignment,
verifiability, steerability, and adaptability rather than on autonomous task
completion alone~\cite{wang2026humansmissing}. The read layer delivers
verifiability (the developer checks what the tool recorded), the write layer
delivers steerability (the developer removes or redirects what was recorded), and
the temporal ordering between them---read first, write later---is adaptability
expressed as deployment strategy rather than runtime architecture.

The temporal strategy is also an acknowledgment about what a designer can know
in advance. Suchman argued that a plan is a resource for situated action rather
than a program that determines it: people use plans to orient themselves and then
improvise in response to what actually happens~\cite{suchman1987}. A design
document that specifies ``ship both layers simultaneously'' is a plan in her
sense, and the evaluation shows what the situated action looked like: the read
layer held on first contact and the write layer did not, in a pattern no amount
of pre-deployment testing predicted at its actual severity. Shipping the read
layer first is the deployment equivalent of a checkpoint: it lets the designer
observe which parts of the plan survived contact with a real user, and defer the
parts that did not until the observation has informed the next design. The
alternative---shipping everything and hoping the plan holds---is what produced
the 50\% per-session violation rate, because the plan said ``composition is
reliable'' and the situated action discovered it was not. A designer who cannot
predict which layer will fail (and our evaluation shows that the failure was in
the write layer, which is the layer whose correctness is harder to verify without
deployment) should ship the layer whose failures are cheapest first and learn
from it.

This recommendation does not assume the read layer is correct. Labels are wrong
on 15\% of features, over-decomposition runs 2--6$\times$ at the median, and the
identity tracker loses a function on roughly half of renames that co-occur with
a rewrite. The read layer is shippable despite these errors because their cost
is bounded: a wrong label costs a rename command, a wrong grouping costs a
regroup, and an identity break costs a manual join---all attention costs, none of
which corrupt data or leave the repository in a state the developer cannot
reverse in one command. A write-layer error of comparable frequency (the 50\%
per-session violation rate) corrupts files, produces silent wrong states, and
erodes the basis for trusting subsequent reports. The asymmetry is what makes the
temporal strategy viable: a read layer whose errors cost attention can earn trust
while being imperfect, whereas a write layer whose errors cost data cannot.

The transfer mechanism is not faith. The pilot participant explicitly
compartmentalised: ``I trust the reports but I don't trust the operations
yet.'' Trust earned through correct attributions did not generalise to write
operations. What transferred was comprehension---a mental model of which
functions belong to which feature, built by reading the tool's reports over
twenty minutes. The cognitive mechanism is that attribution converts a recall
task (``which functions did my request produce?'') into a recognition task
(``does this grouping match what I remember?''), and recognition operates at
lower confidence thresholds than recall for the same
material~\cite{mandler1980}. Twenty minutes of reading attributions builds a
model that twenty minutes of reading code does not, because the attribution
presents the decomposition as something to verify rather than something to
reconstruct. That model is what made the revert preview evaluable: the
participant could judge whether the consequence list was complete because they
already knew what the feature contained. Without that prior model, the same
preview would be an assertion they could not check, and the pilot's corruption
case (a preview that asserted completeness over a result that was wrong) shows
what happens when a developer accepts an uncheckable assertion. The temporal
strategy therefore depends not on trust transfer but on comprehension transfer:
the developer who has read the map can verify the operation's preview
independently, and independent verification is what builds write-layer trust
through demonstrated correctness under observation rather than through faith
inherited from a different mode of use.

Lee and See's review of trust in automation found that a single visible error
erodes trust faster than a string of correct outcomes restores
it~\cite{lee2004}. Trust is harder to earn after a write failure than before
one. This recommendation is in tension with the
paper's own motivation: Section~\ref{sec:scenarios} argues for operations
(revert, restore, switch) that require writing, and a read-only tool cannot
offer them. The resolution is temporal. Ship the read layer now, and ship the
write layer when its failure rate is low enough that the comprehension the read
layer built lets the developer verify each operation's preview before
confirming it. The per-session violation rate of 50\%
(Section~\ref{sec:session-rate}) says we have not crossed that threshold. The
pilot is evidence that we shipped the write layer too early. The read layer also
has a structural role that the trust argument alone does not capture: it is the
external reference that breaks the circularity Zietsman identified in
AI-reviewing-AI pipelines~\cite{zietsman2026}. Without it, the developer's only
check on a revert preview is their own reading of the code---the same reading
that the tool was built to make unnecessary---and the verification degenerates
into the manual recovery the five scenarios describe. Treude's analysis of terms
of service makes the same point from the governance side: providers consistently
shift responsibility for correctness onto the developer while providing no
mechanism for exercising that responsibility at agent
granularity~\cite{treude2026accountable}. The read layer is that mechanism---not
a safety net the provider offers, but an infrastructure the developer builds
for themselves, using the one artefact they do control (the request they typed) as
the anchor for everything the agent subsequently produced.

\subsection{What would change our conclusion}

The result that would overturn this paper's central claim is an error-rate
finding. If a controlled study shows that developers using git alone detect
unintended removals at the same rate as developers using \sgt{}'s preview, then
the comprehension benefit Section~\ref{sec:scenarios} argues for does not exist
in practice, and the fork rule, the layout seam, and the naming instability
provide no benefit. The
pilot cannot settle this. One participant per condition (with the
researcher present) measures tool fitness rather than developer behaviour.
The sgt condition's mutation failures make the comparison one-sided.
The trailer confound (Section~\ref{sec:pilot-trailers}) degrades the git fallback
in the treatment arm, so any comprehension advantage we observed could partly
reflect a weakened baseline rather than a strengthened treatment. The
within-subjects study of Section~\ref{sec:eval-study}, on tasks where both
conditions finish (which the pilot showed they can), answers it in
Section~\ref{sec:findings}. The reach-prediction measure provides the
complementary test: if \texttt{gain} (the movement from blind prediction toward
ground truth after consulting the tool) is zero or negative under \sgt{}, the
representation does not change what the developer expects, and the read layer
claimed more than it delivered. If \texttt{gain} is positive but paired with
high confidence over a wrong answer, the tool confidently misleads---the
inferred-label failure mode this design is most likely to produce.

Two secondary results would reshape specific claims. First, if the fork rule's
amplification factor holds on repositories recorded live (rather than mined
retroactively), then the 0.23 rebuild rate is a property of the design under its
intended workflow, and the distinction Section~\ref{sec:unreproducible} draws
between the two regimes would collapse. We have partial evidence against this: on the two study repositories,
recorded entirely through \cmd{sgt save}, reconstruction is exact (0 drifted
paths), which suggests that the fork rule's cost scales with history length
rather than with recording method. Second, if the codebook's dominant
residual---ground-truth grain disagreement---replicates on repositories with
more intentful request histories, then the pairwise F1 metric is fundamentally
unsuited to this domain and the over-decomposition ratio (2--6$\times$ at the
median, up to 13) is the only stable measurement available. We report both findings so that future work
on intent recovery can calibrate its metrics before discovering the same
instability we found.

A third result would complicate even a positive study finding. Bansal et al.\
showed that AI explanations increase acceptance of recommendations regardless of
whether the recommendation is correct~\cite{bansal2021}. If the study finds
that \sgt{} participants detect more errors than git participants, the mechanism
could be either that the preview \emph{reveals} consequences the participant
would have missed (the mechanism we claim) or that the preview \emph{increases
confidence} in whatever state the participant already believes, including states
that contain undetected errors. The study's third scoring criterion---no function
outside the intended set differs from the reference---distinguishes the two: a
participant who accepts a preview's completeness assertion and submits a
repository with unreported damage has been made \emph{more} wrong by the
explanation, not less. If this pattern appears---participants in the sgt condition
submitting more confidently while carrying hidden damage---then the preview has
the same pathology Bansal et al.\ measured: it makes acceptance feel safer
without making it actually safer. Vasconcelos et al.'s finding that
miscalibrated uncertainty indicators produce worse outcomes than no indicator at
all~\cite{vasconcelos2025} predicts the same failure mode for a preview whose
completeness varies silently across operations. The pre-registered measure
(undetected errors) would catch this, because a participant who trusts a
wrong preview produces more undetected errors, not fewer. We flag this in
advance so that a positive result on the error-rate measure can be read as
genuine rather than as a confidence artefact.

\subsection{The instrument shapes the claim}
\label{sec:instrument-shapes-claim}

Six corrections to published numbers were required during the evaluation, and
each followed the same shape: a rate was reported before the sentence naming its
population was checked against the denominator's actual content. Four of the six
favoured the system. None was a transcription error; all were a denominator whose
\emph{label} implied one set and whose \emph{content} was another. ``Files in a
parsed language'' included 63 Markdown files that reconstruct trivially;
``operations targeting an entity'' included 878 operations that tested a
dependency path no fixture could reach. In each case the arithmetic was correct
and the population sentence was wrong, and the population sentence was what a
reader would take away.

This error has a precise structure that a reader evaluating any self-built system
should recognise. A metric's denominator determines which failures are visible.
A surplus file---one the rebuild invents that does not exist at \cmd{HEAD}---cannot
appear in the numerator of a rate defined over files at \cmd{HEAD}, so a metric
that only asks ``how many of the working tree's files can the record reproduce?''
is structurally blind to the rebuild producing files the working tree never had.
We measured that rate for a month before asking the complementary question, and
the complementary question moved the headline from 0.33 to 0.24. No amount of
careful measurement within a one-sided metric would have found it.

The broader lesson is that the fixture corpus we built to exercise our invariants
was structurally incapable of testing the case real repositories present. Fixtures
round-trip by construction, so the write guard that fires when the store does not
reproduce committed bytes was unreachable. Fixtures write every caller and callee
in a single commit, so the cross-commit dependency path was untestable. Both
constraints followed from reasonable design choices (fixtures should be
deterministic; fixtures should be self-contained), and both made the instruments
blind to the behaviours that dominate on external input. The lesson for any
evaluation of this kind: a corpus designed to \emph{exercise} a system's laws
will systematically avoid the states where those laws have no force, because the
designer optimises for the same conditions the system optimises for. The
evaluator's attention is non-uniform for the same structural reason: a number
that flatters the system looks right, so it passes without the check that would
have found the denominator was wrong. Four of six corrections moved in the same
direction, and each was found by a question we asked about a \emph{different}
number---never by auditing the one we were comfortable with.

A complementary error appeared in the harness itself. Ten thousand operations
reads as thorough, and it is thorough about the guards that check the working
tree against the record. But the subtraction report has four output
branches---the splice report, the broken-reference warning, the carried-forward
notice, and the no-consequence confirmation---and 10{,}237 uniform draws
exercised only two of the four. The two untested branches shared a gate that
required a creation-bearing removal, and creation ops constitute 3--5\% of a
typical live ideal, so a uniform draw over the full population produces too few
creation-bearing removals to reach either branch in any reasonable sample. Volume
measures effort, not coverage: a ten-thousand-operation sweep over one
distribution can be structurally incapable of reaching a code path that a
fifty-operation sweep over a different distribution would exercise on the first
draw. The fix---a weighted draw that over-samples creation ops---found both
branches immediately and produced a defect within a single sitting.

